\slcol{अर्जुन उवाच —\\
\Index{संन्यासं कर्मणां} कृष्ण पुनर्योगं च शंससि ।\\
यच्छ्रेय एतयोरेकं तन्मे ब्रूहि सुनिश्चितम् ॥ ५.१ ॥}
\translate{Arjuna uvāca —\\
saṁnyāsaṁ karmaṇāṁ kṛṣṇa punaryōgaṁ ca śaṁsasi |\\
yacchrēya ētayōrēkaṁ tanmē brūhi suniścitam ‖ 5.1 ‖}
\cquote{Arjuna said:
O Krishna, you extol both the renunciation of action and again, the path of action (\emph{karma-yoga}). Tell me decisively—which of these two is superior?}


\slcol{श्रीभगवानुवाच —\\
\Index{संन्यासः कर्मयोगश्च} निःश्रेयसकरावुभौ ।\\
तयोस्तु कर्मसंन्यासात्कर्मयोगो विशिष्यते ॥ ५.२ ॥}
\translate{śrībhagavānuvāca —\\
saṁnyāsaḥ karmayōgaśca niḥśrēyasakarāvubhau |\\
tayōstu karmasaṁnyāsātkarmayōgō viśiṣyatē ‖ 5.2 ‖}
\cquote{Sri Bhagavan said:\\
Both renunciation and the \emph{yoga} of action (\emph{karma-yoga}) lead to the highest bliss; but of the two, the \emph{yoga} of action is superior (being applicable to all) to the renunciation of action.}

\newpage
\begin{mananam}{\mananamfont {Mananam sloka - 2}} 
\mananamtext Do I take on more activities than is necessary? Am I able to withdraw from activities that are not useful or necessary? On the other hand, do I tend to relinquish my necessary duties and responsibilities, especially when under stress or faced with obstacles? How can I learn to balance what I need to do and what I don’t have to? How can I engage fully in actions while maintaining inner detachment?
\end{mananam}
\WritingHand\enspace{\selfReflect\small{Self Reflection}}
\begin{inspiration}{\mananamfont Inspiration}
\mananamtext Of the two spiritual attitudes towards activity—renunciation and actions without expectations of rewards—the latter offers a greater benefit to the majority who have responsibilities and duties which cannot be given up. However, the wiser way is to know when to push oneself towards action and when to withdraw from activities that are not necessary or helpful on one’s spiritual path. Everyone should cultivate this wisdom, knowing when to swing into action and when to withdraw from it.
\end{inspiration}
\newpage


\slcol{\Index{ज्ञेयः स नित्यसंन्यासी} यो न द्वेष्टि न काङ्क्षति ।\\
निर्द्वन्द्वो हि महाबाहो सुखं बन्धात्प्रमुच्यते ॥ ५.३ ॥}
\translate{jñēyaḥ sa nityasaṁnyāsī yō na dvēṣṭi na kāṅkṣati |\\
nirdvandvō hi mahābāhō sukhaṁ bandhātpramucyatē ‖ 5.3 ‖}
\cquote{One who neither hates nor desires, remaining free from all dualities, should be regarded as an eternal \emph{sannyasi}, O mighty-armed (Arjuna); for such a one is easily freed from the bondage of \emph{samsara}.}

\slcol{\Index{साङ्ख्ययोगौ पृथग्बालाः} प्रवदन्ति न पण्डिताः ।\\
एकमप्यास्थितः सम्यगुभयोर्विन्दते फलम् ॥ ५.४ ॥}
\translate{sāṅ‍khyayōgau pṛthagbālāḥ pravadanti na paṇḍitāḥ |\\
ēkamapyāsthitaḥ samyagubhayōrvindatē phalam ‖ 5.4 ‖}
\cquote{The ignorant perceive the children (the ignorant), not the wise, speak of the path of knowledge and the \emph{yoga} of action as though distinct (in their goals). Truly, being established in either, one attains the fruit of both.}


\slcol{\Index{यत्साङ्ख्यैप्राप्यते स्थानं} तद्योगैरपि गम्यते ।\\
एकं साङ्ख्यं च योगं च यः पश्यति स पश्यति ॥ ५.५ ॥}
\translate{yatsāṅ‍khyaiḥ prāpyatē sthānaṁ tadyōgairapi gamyatē |\\
ēkaṁ sāṅ‍khyaṁ ca yōgaṁ ca yaḥ paśyati sa paśyati ‖ 5.5 ‖}
\cquote{That place which is reached by the \emph{sankhyas} (\emph{jñānīs}) is also reached by the \emph{yogis} (\emph{karma-yogis}). One who sees knowledge and the performance of action (\emph{karma-yoga}) as one has truly perceived.}

\slcol{\Index{संन्यासस्तु महाबाहो} दुःखमाप्तुमयोगतः ।\\
योगयुक्तो मुनिर्ब्रह्म नचिरेणाधिगच्छति ॥ ५.६ ॥}
\translate{saṁnyāsastu mahābāhō duḥkhamāptumayōgataḥ |\\
yōgayuktō munirbrahma nacirēṇādhigacchati ‖ 5.6 ‖}
\cquote{O mighty-armed (Arjuna), renunciation is hard to attain without \emph{karma-yoga}. Thus, one dedicated to this \emph{yoga} quickly attains the Supreme Brahman.}

\newpage
\begin{mananam}{\mananamfont {Mananam sloka - 3}} 
\mananamtext Am I aware of my likes and dislikes and how they condition my choices and decisions in life? Am I able to continue with dutiful actions even when I harbour unconscious dislikes, and discontinue unproductive actions when I harbour a strong like or attachment towards them? How can I learn to function with interest but without attachment?
\end{mananam}
\WritingHand\enspace{\selfReflect\small{Self Reflection}}
\begin{inspiration}{\mananamfont Inspiration}
\mananamtext Inner renunciation is an essential life skill, regardless of whether one is a monk or a householder. Such a one is an eternal \emph{sannyasi} and has perfected the sense of detachment in the mind. He or she may have preferences of a certain lifestyle and habits that are conducive for \emph{sadhana}, but inwardly they harbour neither likes nor dislikes towards anything or anyone. When circumstances warrant it, they are quick to drop anything that they have taken on, and move on with life.
\end{inspiration}
\newpage

\newpage
\begin{mananam}{\mananamfont {Mananam sloka - 6}} 
\mananamtext Do I wish to pursue spiritual life whole heartedly? Am I ready to give up all activities, at least most of them in pursuit of spirituality? Am I qualified to renounce worldly activities or to take up retirement whenever it is due, or even at an earlier juncture? Do I have a spiritual discipline or \emph{sadhana} that I can commit to? How do I see myself engaging my time productively and toward the spiritual direction in which I wish to proceed?
\end{mananam}
\WritingHand\enspace{\selfReflect\small{Self Reflection}}
\begin{inspiration}{\mananamfont Inspiration}
\mananamtext To pursue the spiritual path wholeheartedly, one must commit to a \emph{sadhana}. To renounce worldly life, partly or wholly, without having a disciplined life is a sure cause of sensory indulgence and laziness (otherwise called the \emph{tamasic} life). A traditional spiritual path, consisting of \emph{seva}, worship, meditation and study is the best bet for many who wish to live the spiritual life wholly, regardless of the stage of their life.
\end{inspiration}
\newpage

\slcol{\Index{योगयुक्तो विशुद्धात्मा} विजितात्मा जितेन्द्रियः ।\\
सर्वभूतात्मभूतात्मा कुर्वन्नपि न लिप्यते ॥५. ७ ॥}
\translate{yōgayuktō viśuddhātmā vijitātmā jitēndriyaḥ |\\
sarvabhūtātmabhūtātmā kurvannapi na lipyatē ‖ 5. 7 ‖}
\cquote{One who is dedicated to the \emph{yoga} of action, whose mind is pure, a master of oneself, who has conquered the senses and realised one’s self as the self in all beings, remains untainted by actions—even while performing them.}

\slcol{\Index{नैव किञ्चित्करोमीति} युक्तो मन्येत तत्त्ववित् ।\\
पश्यञ्शृण्वन्स्पृशञ्जिघ्रन्नश्नन्गच्छन्स्वपञ्श्वसन् ॥ ५.८ ॥\\
\Index{प्रलपन् विसृजन्गृह्णन्नु}न्मिषन्निमिषन्नपि ।\\
इन्द्रियाणीन्द्रियार्थेषु वर्तन्त इति धारयन् ॥ ५.९ ॥}
\translate{naiva kiñcitkarōmīti yuktō manyēta tattvavit |\\
paśyañśṛṇvanspṛśañjighrannaśnangacchansvapañśvasan ‖5.8‖\\
pralapan visṛjangṛhṇannunmiṣannimiṣannapi |\\
indriyāṇīndriyārthēṣu vartanta iti dhārayan ‖ 5.9 ‖}
\cquote{The knower of truth, anchored in the self, should contemplate: “I do nothing at all,” even while engaged in activities such as seeing, hearing, touching, smelling, eating, moving, sleeping, breathing, speaking, letting go, grasping, opening and closing the eyes—convinced that the senses function in relation to sense objects.}

\newpage
\begin{mananam}{\mananamfont {Mananam sloka - 8, 9}} 
\mananamtext How can I learn to function in this world without being identified with various actions? Can I create a space between the operations of my sense organs and my detached awareness of them? Or, can I observe that there is a sense of ‘I am’, followed by the subsequent arising sense of ‘I am the doer’? 
Can I cultivate the lived knowledge that what I see is a neutral object until I become the seer who has a positive or negative mental reaction to what I see? Can I observe this in my daily life while noticing how each of these (mere seeing versus prejudiced seeing) leave a different impression on the mind? 
\end{mananam}
\WritingHand\enspace{\selfReflect\small{Self Reflection}}
\begin{inspiration}{\mananamfont Inspiration}
\mananamtext We can all benefit by learning the art of holding onto our self-awareness as we go about our daily activities. It is a means to do all our work and face all our experiences without getting unduly excited or stressed. There is ever a choice within us as to how much we get personally identified and invested in our activities. And even if we do get identified with all of our works, how fast can we withdraw back to our natural peaceful state of beingness?
\end{inspiration}
\newpage


\slcol{\Index{ब्रह्मण्याधाय कर्माणि} सङ्गं त्यक्त्वा करोति यः ।\\
लिप्यते न स पापेन पद्मपत्रमिवाम्भसा ॥ ५.१० ॥}
\translate{brahmaṇyādhāya karmāṇi saṅgaṁ tyaktvā karōti yaḥ |\\
lipyatē na sa pāpēna padmapatramivāmbhasā ‖ 5.10 ‖}
\cquote{One who performs actions by dedicating them to Brahman and relinquishing all attachments is not tainted by sin—just as a lotus petal remains untouched by water.}

\slcol{\Index{कायेन मनसा बुद्ध्या} केवलैरिन्द्रियैरपि ।\\
योगिनः कर्म कुर्वन्ति सङ्गं त्यक्त्वात्मशुद्धये ॥ ५.११ ॥}
\translate{kāyēna manasā buddhyā kēvalairindriyairapi |\\
yōginaḥ karma kurvanti saṅgaṁ tyaktvātmaśuddhayē ‖ 5.11 ‖}
\cquote{\emph{Yogis} engage in actions through body, mind, intellect, and senses—having relinquished attachments—solely for self-purification (from ego)}

\slcol{\Index{युक्तः कर्मफलं त्यक्त्वा} शान्तिमाप्नोति नैष्ठिकीम् ।\\
अयुक्तः कामकारेण फले सक्तो निबध्यते ॥ ५.१२ ॥}
\translate{yuktaḥ karmaphalaṁ tyaktvā śāntimāpnōti naiṣṭhikīm |\\
ayuktaḥ kāmakārēṇa phalē saktō nibadhyatē ‖ 5.12 ‖}
\cquote{One who is dedicated (on the path of \emph{yoga}), having renounced the fruits of action, attains steadfast peace; whereas one lacking dedication, being attached to and driven by the desire for the fruits of action, remains bound.}


\slcol{\Index{सर्वकर्माणि मनसा} संन्यस्यास्ते सुखं वशी ।\\
नवद्वारे पुरे देही नैव कुर्वन्न कारयन् ॥ ५.१३ ॥}
\translate{sarvakarmāṇi manasā saṁnyasyāstē sukhaṁ vaśī |\\
navadvārē purē dēhī naiva kurvanna kārayan ‖ 5.13 ‖}
\cquote{Having mentally renounced all actions, with self-control, the embodied one rests happily in the city of nine gates—neither acting nor causing anyone to act.}

\newpage
\begin{mananam}{\mananamfont {Mananam sloka - 11}} 
\mananamtext What is the nature of the ego in actions? How does it manifest in my life? How can I attain purification of the ego? Is it necessary to renounce all actions that trigger the ego or is there a balanced approach? What does it mean that \emph{yogis} are able to function with detachment? Are they completely free from ego or could there be some kind of a functional ego? What is the nature of the latter?
\end{mananam}
\WritingHand\enspace{\selfReflect\small{Self Reflection}}
\begin{inspiration}{\mananamfont Inspiration}
\mananamtext All the ancient traditions of India emphasized the need for selfless actions as a means of purification of the mind. \emph{Seva} and \emph{karma-yoga} are invaluable for a \emph{sadhak} (spiritual practitioner), regardless of what kind of \emph{sadhana} one embraces. Without it, it is hard to get rid of the ego that functions at different levels and manifests in different ways in a \emph{sadhak’s} life.
\end{inspiration}
\newpage

\begin{mananam}{\mananamfont {Mananam sloka - 13}} 
\mananamtext Am I aware of how the gates of body, such as eyes and ears, are the medium for our interaction with the outside world? Am I aware that what comes in through these gates, and what goes out, as my verbal and physical responses, create an impression on the mind? Am I aware of the moment of choice before my attention and energy goes out through these gates and gets completely identified with the objects of attention?
\end{mananam}
\WritingHand\enspace{\selfReflect\small{Self Reflection}}
\begin{inspiration}{\mananamfont Inspiration}
\mananamtext The inner being within us is the higher self. Its interactions with the outside world causes us to lose our true identity, leading to various challenges in our interactions with others. To counter this, we must cultivate self-awareness or the witnessing consciousness as we engage in daily activities. Start with small steps: observe a flower or a bird in the sky, listen to distant sounds, walk within your home or on your terrace, and pay attention to your own breath. In all these moments, practice maintaining partial awareness of your inner being or the higher self.
\end{inspiration}
\newpage

\slcol{\Index{न कर्तृत्वं न कर्माणि} लोकस्य सृजति प्रभुः ।\\
न कर्मफलसंयोगं स्वभावस्तु प्रवर्तते ॥ ५.१४ ॥}
\translate{na kartṛtvaṁ na karmāṇi lōkasya sṛjati prabhuḥ |\\
na karmaphalasaṁyōgaṁ svabhāvastu pravartatē ‖ 5.14 ‖}
\cquote{The Lord does not create the sense of doership, nor the actions that unite effort with results; it is one’s own nature that manifests outwardly.}
\slcol{\Index{नादत्ते कस्यचित्पापं} न चैव सुकृतं विभुः ।\\
अज्ञानेनावृतं ज्ञानं तेन मुह्यन्ति जन्तवः ॥ ५.१५ ॥} 
\translate{nādattē kasyacitpāpaṁ na caiva sukṛtaṁ vibhuḥ |\\
ajñānēnāvṛtaṁ jñānaṁ tēna muhyanti jantavaḥ ‖ 5.15 ‖}
\cquote{The Lord does not dispense anyone’s virtue or sin. Knowledge is veiled by one’s own ignorance, and thereby beings are deluded.}

\slcol{\Index{ज्ञानेन तु तदज्ञानं} येषां नाशितमात्मनः ।\\
तेषामादित्यवज्ज्ञानं प्रकाशयति तत्परम् ॥ ५.१६ ॥}
\translate{jñānēna tu tadajñānaṁ yēṣāṁ nāśitamātmanaḥ |\\
tēṣāmādityavajjñānaṁ prakāśayati tatparam ‖ 5.16 ‖}
\cquote{But for those whose ignorance is destroyed by the knowledge of the self—like the sun revealing the world—such knowledge reveals the Supreme.}

\slcol{\Index{तद्बुद्धयस्तदात्मान}स्तन्निष्ठास्तत्परायणाः ।\\
गच्छन्त्यपुनरावृत्तिं ज्ञाननिर्धूतकल्मषाः ॥५.१७ ॥} 
\translate{tadbuddhayastadātmānastanniṣṭhāstatparāyaṇāḥ |\\
gacchantyapunarāvṛttiṁ jñānanirdhūtakalmaṣāḥ ‖ 5.17 ‖}
\cquote{Those whose intellect is absorbed in That, whose self is That, who are established in That, with That as their supreme goal—attain the state of non-returning, their sins dispelled by knowledge.}

\newpage
\begin{mananam}{\mananamfont {Mananam sloka - 16, 17}} 
\mananamtext Why do the ancient scriptures tell us we are ignorant? What areas of knowledge elude me despite my academic education and civilised upbringing? Am I capable of discerning the distinction between worldly knowledge and spiritual wisdom? Which knowledge can lead me to the supreme truth? What practices can help align my mind with that ultimate reality? In what ways can these practices transform my daily life? 
\end{mananam}
\WritingHand\enspace{\selfReflect\small{Self Reflection}}
\begin{inspiration}{\mananamfont Inspiration}
\mananamtext The culmination of all spiritual paths is the eradication of ignorance. Just as darkness yields to light, knowledge alone can dispel ignorance. One insidious effect of ignorance is its concealment of our true identity. Consequently, all our actions arise from this darkness. Continued efforts to establish our intellect in the highest wisdom breaks down all limitations.
\end{inspiration}
\newpage


\slcol{\Index{विद्याविनयसम्पन्ने} ब्राह्मणे गवि हस्तिनि ।\\
शुनि चैव श्वपाके च पण्डिताः समदर्शिनः ॥ ५.१८ ॥}
\translate{vidyāvinayasampannē brāhmaṇē gavi hastini |\\
śuni caiva śvapākē ca paṇḍitāḥ samadarśinaḥ ‖ 5.18 ‖}
\cquote{The sage, endowed with learning and humility, looks upon a Brahmana, a cow, an elephant, a dog, and even an eater of dog’s meat (a low-born one) with equal vision—without discrimination.}

\slcol{\Index{इहैव तैर्जितः सर्गो} येषां साम्ये स्थितं मनः ।\\
निर्दोषं हि समं ब्रह्म तस्माद्ब्रह्मणि ते स्थिताः ॥ ५.१९ ॥} 
\translate{ihaiva tairjitaḥ sargō yēṣāṁ sāmyē sthitaṁ manaḥ |\\
nirdōṣaṁ hi samaṁ brahma tasmādbrahmaṇi tē sthitāḥ ‖ 5.19 ‖}
\cquote{Here in this world, those whose mind remains established in equanimity overcome birth and death. Since Brahman is flawless and impartial in all, they are established in that Brahman.}
\slcol{\Index{न प्रहृष्येत्प्रियं} प्राप्य नोद्विजेत्प्राप्य चाप्रियम् ।\\
स्थिरबुद्धिरसंमूढो ब्रह्मविद्ब्रह्मणि स्थितः ॥ ५.२० ॥}
\translate{na prahṛṣyētpriyaṁ prāpya nōdvijētprāpya cāpriyam |\\
sthirabuddhirasaṁmūḍhō brahmavidbrahmaṇi sthitaḥ ‖ 5.20 ‖}
\cquote{A knower of Brahman, with steady intellect and undeluded, is neither elated by the pleasant nor dejected by the unpleasant; such a one is established in Brahman.}
\slcol{\Index{बाह्यस्पर्शेष्वसक्तात्मा} विन्दत्यात्मनि यत्सुखम् ।\\
स ब्रह्मयोगयुक्तात्मा सुखमक्षयमश्नुते ॥ ५.२१ ॥}
\translate{bāhyasparśēṣvasaktātmā vindatyātmani yatsukham |\\
sa brahmayōgayuktātmā sukhamakṣayamaśnutē ‖ 5.21 ‖}
\cquote{With the mind detached from external objects, one finds joy in the self; with the mind absorbed in meditation on Brahman, one attains everlasting bliss.}


\newpage
\begin{mananam}{\mananamfont {Mananam sloka - 18}} 
\mananamtext By what criteria do I evaluate people in society – based on their status, riches, education, beauty, or achievements? Whom do I honour more and whom do I honour less? Do I hold some people in lower regard? How can I benefit by giving my attention and care to those who are needy? Who are the role models that I can look up to, those whose values and life experiences may elevate my intellectual and spiritual growth?
\end{mananam}
\WritingHand\enspace{\selfReflect\small{Self Reflection}}
\begin{inspiration}{\mananamfont Inspiration}
\mananamtext An exalted human being is one whose response remains equanimous toward all, neither awed by the attention of a rich, powerful individual nor disdainful of someone with low social status.
Equanimity should not be mistaken for indifference. Instead, we must learn to recognize those who require our attention and those who deserve reverence. As a practice in shaping our value system, we must shift our focus from our impulse of honouring the wealthy and powerful to prioritising our respect for those who are spiritually wise. Such a value system is a sign of a spiritually advanced society.
\end{inspiration}
\newpage



\newpage
\begin{mananam}{\mananamfont {Mananam sloka - 20}} 
\mananamtext When faced with pleasant experiences, how do I respond? How do I handle unpleasant situations? Am I convinced that rising above the seeking of pleasure and the avoidance of pain is beneficial for my mental well-being? If so, how?
Am I conscious of the choices between immediate pleasure and the long-term good in all aspects of life? What is the value of the calm acceptance of results or outcomes, especially after I have put in significant efforts?
\end{mananam}
\WritingHand\enspace{\selfReflect\small{Self Reflection}}
\begin{inspiration}{\mananamfont Inspiration}
\mananamtext What is often seen as a ‘natural impulse’ for humans is regarded as harmful by wise sages and \emph{yogis}. Almost everyone rushes to seek what is pleasant and avoid what is unpleasant. However, it is precisely these likes and dislikes that lead to inner instability. Cultivating mental equanimity by accepting both pleasant and unpleasant experiences is a highly beneficial practice.
\end{inspiration}
\newpage


\slcol{\Index{ये हि संस्पर्शजा भोगा} दुःखयोनय एव ते ।\\
आद्यन्तवन्तः कौन्तेय न तेषु रमते बुधः ॥ ५.२२ ॥}
\translate{yē hi saṁsparśajā bhōgā duḥkhayōnaya ēva tē |\\
ādyantavantaḥ kauntēya na tēṣu ramatē budhaḥ ‖ 5.22 ‖}
\cquote{Pleasurable experiences born of contact with objects have a beginning and an end, O son of Kunti. The wise do not rejoice in them (nor do they yearn for them).}

\slcol{\Index{शक्नोतीहैव यः सोढुं} प्राक्छरीरविमोक्षणात् ।\\
कामक्रोधोद्भवं वेगं स युक्तः स सुखी नरः ॥ ५.२३ ॥}
\translate{śaknōtīhaiva yaḥ sōḍhuṁ prākcharīravimōkṣaṇāt |\\
kāmakrōdhōdbhavaṁ vēgaṁ sa yuktaḥ sa sukhī naraḥ ‖ 5.23 ‖}
\cquote{One who is able, here in this world before giving up the body, to remain resilient against the impulses of desire and anger—such a one is a \emph{yogi} and a happy being.}


\slcol{\Index{योऽन्तःसुखोऽन्तराराम}स्तथान्तर्ज्योतिरेव यः ।\\
स योगी ब्रह्मनिर्वाणं ब्रह्मभूतोऽधिगच्छति ॥ ५.२४ ॥}
\translate{yō:'ntaḥsukhō:'ntarārāmastathāntarjyōtirēva yaḥ |\\
sa yōgī brahmanirvāṇaṁ brahmabhūtō:'dhigacchati ‖ 5.24 ‖}
\cquote{One who finds joy within, delights within, and is illumined within—such a \emph{yogi}, having become Brahman, attains liberation in Brahman.}


\slcol{\Index{लभन्ते ब्रह्मनिर्वाणमृषयः} क्षीणकल्मषाः ।\\ 
छिन्नद्वैधा यतात्मानः सर्वभूतहिते रताः ॥ ५.२५ ॥}
\translate{labhantē brahmanirvāṇamṛṣayaḥ kṣīṇakalmaṣāḥ |\\
chinnadvaidhā yatātmānaḥ sarvabhūtahitē ratāḥ ‖ 5.25 ‖}
\cquote{Those sages whose sins are destroyed, for whom the pairs of opposites are eliminated, who are self-controlled and devoted to the welfare of all beings—such ones attain liberation in Brahman.}

\newpage
\begin{mananam}{\mananamfont {Mananam sloka - 21, 22}} 
\mananamtext Am I conscious of the interaction between my senses and sense-objects? For instance, when my eyes perceive something attractive, am I aware of the mental reactions it triggers—such as the desire to acquire or to possess?
Do external objects—like food, entertainment, or social interaction—truly bring lasting happiness? Or are they only capable of providing short-term satisfaction? Am I aware of the impermanence inherent in seeking happiness solely from external sources?
\end{mananam}
\WritingHand\enspace{\selfReflect\small{Self Reflection}}
\begin{inspiration}{\mananamfont Inspiration}
\mananamtext We often seek happiness in impermanent external objects. Consequently, this happiness, too, becomes fleeting and leads to suffering when those objects vanish. Our inner thoughts and emotions are closely tied to the external world. Happiness and sorrow are experienced within our minds. 
Seeing this clearly gives one the power to disconnect the outer world from the inner world. To learn to be happy at will, regardless of outer circumstances, is an essential skill on the spiritual path.
\end{inspiration}
\newpage


\slcol{\Index{कामक्रोधवियुक्तानां} यतीनां यतचेतसाम् ।\\
अभितो ब्रह्मनिर्वाणं वर्तते विदितात्मनाम् ॥ ५.२६ ॥} 
\translate{kāmakrōdhaviyuktānāṁ yatīnāṁ yatacētasām |\\
abhitō brahmanirvāṇaṁ vartatē viditātmanām ‖ 5.26 ‖}
\cquote{Those self-controlled ones, free from desire and anger, with minds well-governed and the self realised—attain liberation in Brahman, whether in life or after death.}

\slcol{\Index{स्पर्शान्कृत्वा} बहिर्बाह्यांश्चक्षुश्चैवान्तरे भ्रुवोः ।\\
प्राणापानौ समौ कृत्वा नासाभ्यन्तरचारिणौ ॥ ५.२७ ॥}
\translate{sparśānkṛtvā bahirbāhyāṁścakṣuścaivāntarē bhruvōḥ |\\
prāṇāpānau samau kṛtvā nāsābhyantaracāriṇau ‖ 5.27 ‖}
\cquote{Keeping out all external objects and fixing the eyes at the point between the eyebrows, equalising the outgoing and incoming breaths moving through the nostrils -}
\slcol{\Index{यतेन्द्रियमनोबुद्धि}र्मुनिर्मोक्षपरायणः ।\\ 
विगतेच्छाभयक्रोधो यः सदा मुक्त एव सः ॥ ५.२८ ॥}
\translate{yatēndriyamanōbuddhirmunirmōkṣaparāyaṇaḥ |\\ 
vigatēcchābhayakrōdhō yaḥ sadā mukta ēva saḥ ‖ 5.28 ‖} 
\cquote{The sage, with senses, mind, and intellect restrained, having liberation as the supreme goal, free from desire, fear, and anger—such a one is liberated forever.}
\slcol{\Index{भोक्तारं यज्ञतपसां} सर्वलोकमहेश्वरम् ।\\ 
सुहृदं सर्वभूतानां ज्ञात्वा मां शान्तिमृच्छति ॥ ५.२९ ॥ }
\translate{bhōktāraṁ yajñatapasāṁ sarvalōkamahēśvaram |\\
suhṛdaṁ sarvabhūtānāṁ jñātvā māṁ śāntimṛcchati ‖ 5.29 ‖} 
\cquote{One attains supreme peace who knows Me as the enjoyer of all sacrifices and austerities, the great Lord of all worlds, and the friend of all beings.}

\newpage
\begin{mananam}{\mananamfont {Mananam sloka - 26}} 
\mananamtext What role do desires and related emotions play in my life? Are my desires healthy and constructive for my long-term well-being? Is my motivation in life driven solely by self-centred desires? Can I find motivation to act based on the intention of the well-being of others? Can I recall instances where strong attachment led to anger due to perceived obstruction of my desires? What does it mean to me to attain the state of the sages, devoid of all passions and cravings? 
\end{mananam}
\WritingHand\enspace{\selfReflect\small{Self Reflection}}
\begin{inspiration}{\mananamfont Inspiration}
\mananamtext The common person’s concept of freedom involves having all his needs and desires fulfilled as soon as they arise in his mind. In contrast, \emph{yogis} attain freedom through vigilant mental awareness, where they allow no unwanted desires for sensory gratification or egoic validation. If any desires do arise, they are immediately transmuted and directed toward the well-being of others. The very act of desiring ties us down with our body identification. True freedom is to abide in the soul’s natural state of supreme bliss.
\end{inspiration}
\newpage

\begin{center}
{\devfont  ॐ तत्सदिति श्रीमद्भगवद्गीतासूपनिषत्सु ब्रह्मविद्यायां योगशास्त्रे\\
  श्रीकृष्णार्जुनसम्वादे संन्यासयोगो नाम पञ्चमोऽध्यायः॥\\}
  \chapEndSlokaTranslate{ōṃ tatsaditi śrīmadbhagavadgītās ūpaniṣatsu \\brahmavidyāyāṃ yōgaśāstrē\\ śrīkṛṣṇārjunasaṃvādē sannyāsa-yoga\\nāma panchamodhyāyaḥ‖}
\end{center}